\section{端到端实践：怎样审计一条数据管线}

把本讲落到工程流程，可以使用以下检查表：

\begin{longtable}{L{0.17\textwidth}L{0.35\textwidth}L{0.38\textwidth}}
\toprule
层次 & 必须记录 & 需要抽查的问题 \\
\midrule
Source & URL、license、crawl time、content type & 是否允许使用，是否存在地域/语言偏差 \\
Transform & parser/OCR/version、失败码、结构字段 & 表格、公式、代码顺序是否被破坏 \\
Filter & target data、features、score、threshold & 哪些 domain 被误杀，是否学到 style shortcut \\
Dedup & normalization、粒度、hash/signature、threshold & 是否仍有 contamination，是否删除有用模板 \\
Mix & domain weights、token counts、sampling temperature & proxy ordering 是否能外推到目标规模 \\
Synthetic & generator、prompt、sampling、verifier、tests & 是否存在 judge hacking、自我强化错误 \\
Evaluation & held-out data、freshness、slice metrics & 是否只优化公开 benchmark \\
\bottomrule
\end{longtable}

\begin{importantbox}{最终结论}
强模型的数据优势往往不来自某个神秘数据集，而来自一条可持续迭代的管线：能看到真实样本，能定位错误来源，能通过小实验调整策略，能把 provenance 和 evaluation 接回训练决策。
\end{importantbox}

